%% kapa: exam=1 round=29 year=2018 period=1 form=A subject=real_estate subject_ko=부동산학원론 seq=2 number=82
%% source: Q-Net 감정평가사 기출문제 (cst00309) artlSeq=5202808 file=extracted/1008/5202808/2018년 제29회 감정평가사 1차시험 기출문제/2018년 제29회 감정평가사 1차시험 기출문제/1교시 A형.pdf
%% conversion: kordoc 4.24.51 -> tools/gen1.py (auto); review=proofread-sonnet; flags=none
\begin{question}{82}
부동산활동과 관련된 다음의 내용을 설명하는 용어로 옳게 연결된 것은?
\begin{framedtext}
ㄱ. 인근지역의 주위환경 등의 사정으로 보아 현재의 용도에서 장래 택지 등 다른 용도로의 전환이 객관적으로 예상되는 토지 \par
ㄴ. 택지 등 다른 용도로 조성되기 이전 상태의 토지 \par
\end{framedtext}
\begin{choices}
  \item ㄱ: 후보지, ㄴ: 소지
  \item ㄱ: 후보지, ㄴ: 공지
  \item ㄱ: 이행지, ㄴ: 소지
  \item ㄱ: 이행지, ㄴ: 공지
  \item ㄱ: 성숙지, ㄴ: 소지
\end{choices}
\end{question}
\answer{1}
